\section{Results}
\label{sec:results}

\subsection{CPU results (Apple M1 Pro)}
Deposit-kernel runtimes, bandwidth, charge error, energy drift, OpenMP scaling, and layout/sort breakdown follow the archived \texttt{benchmark.csv} (Figs.~\ref{fig:runtime}--\ref{fig:sort}).
At $N_p=10^5$ on a $128^2$ grid with SoA sorted layout, CIC completes in 0.59\,ms, TSC in 0.90\,ms, and Esirkepov in 4.33\,ms; TSC achieves the lowest spectral noise (0.26).
Eight-thread OpenMP privatization reduces CIC deposit time to 0.18\,ms at the same settings.

\begin{figure}[t]
    \centering
    \includegraphics[width=0.85\linewidth]{fig1_runtime_vs_particles.pdf}
    \caption{Deposit-kernel runtime vs.\ $N_p$ (M1 Pro, $128^2$, SoA, sorted).}
    \label{fig:runtime}
\end{figure}

\subsection{Full-timestep breakdown and amortized sorting}
Figure~\ref{fig:timestep} reports average per-stage times from \texttt{pic\_benchmark --timestep} for CIC with SoA sorting at $128^2$.
Poisson FFT and cell sorting dominate at moderate $N_p$ when sorting runs every step; deposition remains sub-millisecond at $N_p=10^5$.
Figure~\ref{fig:amortized} shows that sorting every $K=10$ or $K=100$ steps reduces amortized sort cost proportionally, making deposit and Poisson competitive---a realistic production-PIC assumption.
Representative deposit times include run-to-run spread (Table~\ref{tab:summary}, mean $\pm$ std over five repeats).

\begin{figure}[t]
    \centering
    \includegraphics[width=0.85\linewidth]{fig13_timestep_breakdown.pdf}
    \caption{Full PIC timestep breakdown (CIC, SoA, sorted every step, M1 Pro).}
    \label{fig:timestep}
\end{figure}

\begin{figure}[t]
    \centering
    \includegraphics[width=0.85\linewidth]{fig15_amortized_sort.pdf}
    \caption{Amortized sort interval $K$ vs.\ per-timestep cost ($N_p=10^5$, $128^2$, CIC).}
    \label{fig:amortized}
\end{figure}

\begin{table}[t]
\centering
\small
\caption{Representative CPU deposition results ($N_p=10^5$, $128^2$ grid, SoA, sorted).}
\label{tab:summary}
\begin{tabular}{lrrrr}
\toprule
Scheme & Sort & Deposit & Thrpt. & $\epsilon_Q$ \\
 & (ms) & (ms) & ($10^8$ p/s) & \\
\midrule
NGP & 2.267 & 1.910 $\pm$ 0.043 & 0.52 & 0.00e+00 \\
CIC & 2.263 & 0.607 $\pm$ 0.061 & 1.65 & 1.16e-17 \\
TSC & 2.421 & 0.848 $\pm$ 0.003 & 1.18 & 1.36e-18 \\
Esirkepov & 2.617 & 7.263 $\pm$ 1.598 & 0.14 & 1.03e-03 \\
\bottomrule
\end{tabular}
\end{table}


\subsection{GPU results (Tesla T4, same host)}
GPU atomic deposition achieves $5.3\times$ speedup over T4-host CPU for CIC at $N_p=10^5$ (4.24\,ms $\rightarrow$ 0.81\,ms); see Figs.~\ref{fig:gpu_runtime} and~\ref{fig:gpu_speedup}.
Multi-block privatized tiles preserve charge conservation but remain $\sim$6$\times$ slower than atomics because each tile scans all particles (Table~\ref{tab:gpu_priv}, Appendix~\ref{app:gpu}).

\begin{figure}[t]
    \centering
    \includegraphics[width=0.85\linewidth]{fig8_gpu_runtime.pdf}
    \caption{CPU vs.\ GPU atomic runtime (T4, $128^2$).}
    \label{fig:gpu_runtime}
\end{figure}

\subsection{Roofline and cross-platform analysis}
Figure~\ref{fig:roofline} places deposition in the memory-bandwidth-limited region; Figure~\ref{fig:bandwidth} reports effective GB/s vs.\ grid size.
Table~\ref{tab:cross_platform} in Section~\ref{sec:code_comparison} summarizes labeled M1 vs.\ T4 metrics.

\begin{figure}[t]
    \centering
    \includegraphics[width=0.85\linewidth]{fig10_roofline.pdf}
    \caption{Roofline-style deposition performance ($N_p=10^5$, $128^2$).}
    \label{fig:roofline}
\end{figure}

\begin{figure}[t]
    \centering
    \includegraphics[width=0.85\linewidth]{fig2_bandwidth_vs_grid.pdf}
    \caption{Effective memory bandwidth vs.\ grid size ($10^5$ particles, M1 Pro).}
    \label{fig:bandwidth}
\end{figure}

\begin{table}[t]
\centering
\small
\caption{Cross-platform CIC deposition ($128^2$ grid, SoA).}
\label{tab:cross_platform}
\begin{tabular}{lrrrr}
\toprule
$N_p$ & M1 CPU (ms) & M1 8T (ms) & T4 CPU (ms) & T4 GPU (ms) \\
\midrule
1e+05 & 0.57 & nan & 4.24 & 0.81 \\
1e+06 & 6.04 & nan & 43.98 & 6.07 \\
\bottomrule
\end{tabular}
\end{table}


\begin{table}[t]
\centering
\small
\caption{CPU vs GPU atomic deposition on Tesla T4 ($N_p=10^5$, $128^2$ grid, SoA).}
\label{tab:gpu_summary}
\begin{tabular}{lrrr}
\toprule
Scheme & CPU (ms) & GPU (ms) & Speedup \\
\midrule
NGP & 3.55 & 0.80 & 4.4$\times$ \\
CIC & 4.24 & 0.81 & 5.3$\times$ \\
TSC & 5.46 & 0.86 & 6.3$\times$ \\
\bottomrule
\end{tabular}
\end{table}

\begin{table}[t]
\centering
\small
\caption{GPU atomic vs privatized CIC on Tesla T4 ($128^2$ grid, SoA).}
\label{tab:gpu_priv}
\begin{tabular}{lrrrr}
\toprule
$N_p$ & Atomics (ms) & Priv (ms) & Atomics speedup & Priv speedup \\
\midrule
1e+04 & 0.18 & 0.57 & 2.3$\times$ & 0.7$\times$ \\
1e+05 & 0.81 & 4.72 & 5.3$\times$ & 0.9$\times$ \\
1e+06 & 6.07 & 39.59 & 7.3$\times$ & 1.1$\times$ \\
\bottomrule
\end{tabular}
\end{table}

